\documentclass[12pt, a4paper]{article}

% ─── Packages ─────────────────────────────────────────────
\usepackage[margin=1in]{geometry}
\usepackage{hyperref}
\usepackage{enumitem}
\usepackage{xcolor}
\usepackage{tcolorbox}
\usepackage{fancyhdr}
\usepackage{titlesec}
\usepackage{parskip}
\usepackage{listings}
\usepackage{graphicx}

% ─── Colors ───────────────────────────────────────────────
\definecolor{promptbg}{RGB}{240, 240, 255}
\definecolor{promptborder}{RGB}{100, 100, 240}
\definecolor{reasonbg}{RGB}{240, 255, 240}
\definecolor{reasonborder}{RGB}{60, 180, 80}
\definecolor{impactbg}{RGB}{255, 248, 230}
\definecolor{impactborder}{RGB}{220, 160, 50}
\definecolor{darkblue}{RGB}{30, 30, 100}

% ─── Custom Boxes ─────────────────────────────────────────
\newtcolorbox{promptbox}[1][]{
  colback=promptbg, colframe=promptborder,
  fonttitle=\bfseries, title={\textcolor{white}{Prompt}}, #1
}
\newtcolorbox{reasonbox}[1][]{
  colback=reasonbg, colframe=reasonborder,
  fonttitle=\bfseries, title={\textcolor{white}{Reasoning \& Design Rationale}}, #1
}
\newtcolorbox{impactbox}[1][]{
  colback=impactbg, colframe=impactborder,
  fonttitle=\bfseries, title={\textcolor{white}{Impact on System}}, #1
}

% ─── Header / Footer ─────────────────────────────────────
\pagestyle{fancy}
\fancyhf{}
\fancyhead[L]{\small Nexus AI — AI Prompts \& Reasoning}
\fancyhead[R]{\small Muhammad Ahmad Abbas}
\fancyfoot[C]{\thepage}

% ─── Title Formatting ────────────────────────────────────
\titleformat{\section}{\Large\bfseries\color{darkblue}}{}{0em}{}[\titlerule]
\titleformat{\subsection}{\large\bfseries\color{darkblue}}{}{0em}{}

% ─── Document ─────────────────────────────────────────────
\begin{document}

\begin{center}
    {\LARGE\bfseries AI Prompts \& Reasoning Document}\\[6pt]
    {\large Nexus AI — Multi-Agent Research Assistant}\\[4pt]
    {\normalsize Muhammad Ahmad Abbas}\\[2pt]
    {\small \today}
\end{center}

\vspace{10pt}

\section*{Executive Summary}

This document details the AI-assisted development process used to build the Nexus AI Multi-Agent Research Assistant. Each section presents a specific prompt I used, the \textbf{reasoning} behind why I crafted it that way, and the \textbf{design impact} it had on the final system. My approach was guided by three core principles:

\begin{enumerate}[leftmargin=2em]
    \item \textbf{Architecture-First Thinking} — I started by defining the system structure, not writing code.
    \item \textbf{Incremental, Test-Driven Development} — Each feature was built, tested, and debugged before moving on.
    \item \textbf{Production-Grade Standards} — I deliberately prompted for industry patterns (low coupling, high cohesion, design patterns) rather than accepting a quick-and-dirty prototype.
\end{enumerate}

\newpage

%% ═══════════════════════════════════════════════════════════
\section{Prompt 1 — Task Decomposition \& Architecture Planning}
%% ═══════════════════════════════════════════════════════════

\begin{promptbox}
\textit{``I have to complete this project within a tight deadline but I have to do it quick and perfect — breakdown the task.''}
\end{promptbox}

\begin{reasonbox}
Before writing a single line of code, I started by requesting a \textbf{full task decomposition}. This is a deliberate application of the software engineering principle of \textbf{planning before implementation}. Rather than diving into code, I wanted:

\begin{itemize}[leftmargin=1.5em]
    \item A clear mental model of all four agents and their responsibilities.
    \item A defined state schema before any node logic was written.
    \item An explicit routing map (Clarity $\to$ Research $\to$ Validator $\to$ Synthesis) with conditional edges.
    \item A realistic timeline to ensure I could deliver a polished, complete product — not an unfinished prototype.
\end{itemize}

\textbf{Design Principle Applied:} \textit{Separation of Concerns} — By breaking the system into discrete agents with single responsibilities, the architecture naturally emerged as a modular pipeline.
\end{reasonbox}

\begin{impactbox}
This prompt produced the foundational \texttt{implementation\_plan.md} artifact that guided the entire project. It established the \texttt{AgentState} schema, the 4-node graph topology, and the conditional routing rules — all before any Python file was created. This prevented scope creep and ensured every subsequent prompt was building toward a clearly defined architecture.
\end{impactbox}

%% ═══════════════════════════════════════════════════════════
\section{Prompt 2 — Technology Selection \& API Strategy}
%% ═══════════════════════════════════════════════════════════

\begin{promptbox}
\textit{``Don't do anything for now — just tell me which free API key is best for my task.''}
\end{promptbox}

\begin{reasonbox}
I deliberately paused implementation to evaluate the \textbf{technology stack} before committing to it. This reflects a critical engineering skill: \textbf{technology due diligence}. Key considerations:

\begin{itemize}[leftmargin=1.5em]
    \item \textbf{Tavily over SerpAPI/Google Custom Search:} Tavily is designed for AI agents — it returns clean, parsed data rather than raw HTML, reducing the preprocessing burden.
    \item \textbf{Groq over OpenAI/Anthropic:} With 4 agents in a loop, latency compounds. Groq's inference speed makes the app feel responsive. Its free tier also removes cost as a barrier.
    \item \textbf{FastAPI over Flask:} I chose FastAPI for its async support, automatic OpenAPI docs, and Pydantic integration — all production-grade features that align with the assignment's quality expectations.
\end{itemize}

\textbf{Design Principle Applied:} \textit{Right Tool for the Right Job} — Each technology was selected based on the specific constraints of the project (speed, cost, agent compatibility), not just familiarity.
\end{reasonbox}

\begin{impactbox}
This evaluation prevented a critical production bug: the original model (\texttt{llama3-8b-8192}) had been decommissioned by Groq. Because I approached technology selection as a deliberate decision point, I was able to quickly diagnose and fix this when it surfaced, switching to \texttt{llama-3.1-8b-instant} without any architectural changes.
\end{impactbox}

%% ═══════════════════════════════════════════════════════════
\section{Prompt 3 — Code Quality \& Project Structure}
%% ═══════════════════════════════════════════════════════════

\begin{promptbox}
\textit{``Let's maintain a README.md file which we will update at each increment. Also the code should be very clean and industry grade — well documented, well organized, and a good directory structure with everything well organized.''}
\end{promptbox}

\begin{reasonbox}
This prompt explicitly prioritizes \textbf{software craftsmanship} over speed. I requested:

\begin{itemize}[leftmargin=1.5em]
    \item \textbf{Incremental README tracking:} A living document that evolves with the codebase — mirrors real-world Agile documentation practices.
    \item \textbf{Modular directory structure:} Separating \texttt{src/agent/}, \texttt{src/api/}, and \texttt{src/web/} enforces the \textbf{Layered Architecture} pattern, ensuring the agent logic has zero coupling to the HTTP layer.
    \item \textbf{Docstrings and comments:} Every function, module, and class is documented — not with ``what'' comments, but ``why'' comments.
\end{itemize}

\textbf{Design Pattern Applied:} \textit{Layered Architecture} — Agent Layer (domain logic) $\to$ API Layer (application logic) $\to$ Presentation Layer (UI). Each layer only depends on the one below it.
\end{reasonbox}

\begin{impactbox}
The directory structure that resulted:
\begin{verbatim}
src/
├── agent/   (state.py, nodes.py, graph.py, utils.py)
├── api/     (main.py)
├── web/     (index.html, styles.css, app.js)
└── config.py
\end{verbatim}
This made every subsequent change (adding \texttt{utils.py}, adding \texttt{config.py}) trivial because the architecture had clear seams for extension — the \textbf{Open/Closed Principle} in action.
\end{impactbox}

%% ═══════════════════════════════════════════════════════════
\section{Prompt 4 — Systematic Testing Strategy}
%% ═══════════════════════════════════════════════════════════

\begin{promptbox}
\textit{``Give me some test cases on which I test it — don't do it by yourself, I'll test it myself. You just provide me with test cases to test it better.''}
\end{promptbox}

\begin{reasonbox}
Rather than trusting that the code worked, I explicitly requested a \textbf{structured test plan} organized by feature area:

\begin{itemize}[leftmargin=1.5em]
    \item \textbf{Group 1 — Clarity Agent:} Ambiguous queries, clear queries, edge cases.
    \item \textbf{Group 2 — Research Pipeline:} Well-known companies, niche entities.
    \item \textbf{Group 3 — Multi-Turn Memory:} Sequential follow-up questions using pronouns.
    \item \textbf{Group 4 — Edge Cases:} Empty input, gibberish, rapid-fire messages.
    \item \textbf{Group 5 — Session Reset:} Verifying state isolation after page refresh.
\end{itemize}

I chose to run these tests \textbf{manually myself} rather than having the AI run them — this ensures I understand the system's behavior firsthand and can catch nuances an automated check might miss.

\textbf{Design Principle Applied:} \textit{Test-Driven Verification} — Tests were designed before bugs were found, not in reaction to them.
\end{reasonbox}

\begin{impactbox}
This test plan immediately uncovered \textbf{two critical bugs}:
\begin{enumerate}
    \item The Clarity Agent didn't use conversation history, causing follow-up questions (``What about their competitors?'') to be incorrectly flagged as ambiguous.
    \item The Synthesis Agent exceeded Groq's 6,000 TPM free-tier limit when processing large Tavily results.
\end{enumerate}
Both were fixed before the final release — without this systematic testing approach, they would have shipped as production bugs.
\end{impactbox}

%% ═══════════════════════════════════════════════════════════
\section{Prompt 5 — Advanced System Design Improvements}
%% ═══════════════════════════════════════════════════════════

\begin{promptbox}
\textit{``Optimize Tavily Search Results. Implement Smart Conversation History using the Sliding Window + Summary Method. Use different models on the free tier for different tasks — like for query refinement, summarizing search results, and finalizing research synthesis. Use models according to our need and requirement and it should work on high accuracy by utilizing the models in an efficient manner. Don't implement any of the above — suggest me more to improve our system design.''}
\end{promptbox}

\begin{reasonbox}
This is arguably the most important prompt in the entire session. Instead of asking for a fix, I proposed \textbf{multiple architectural improvements} and then asked the AI to \textbf{critique and extend} my thinking. This demonstrates:

\begin{itemize}[leftmargin=1.5em]
    \item \textbf{Multi-Model Strategy:} I recognized that not all agents need the same LLM capability. Classification tasks (Clarity, Validator) can use a lightweight 8B model, while the user-facing Synthesis Agent benefits from the 70B model. This is a real-world \textbf{cost-optimization pattern} used in production AI systems.
    \item \textbf{Sliding Window + Summary Memory:} Instead of naive truncation, I proposed keeping recent messages verbatim and summarizing older context — a technique inspired by research in long-context conversation management.
    \item \textbf{Tavily Optimization:} I identified that we were fetching too many results and not leveraging Tavily's \texttt{include\_answer} feature, which provides a pre-summarized direct answer.
    \item \textbf{``Don't implement — suggest more'':} I deliberately asked for \textit{additional ideas before implementation}, ensuring the design was comprehensive before committing to code.
\end{itemize}

\textbf{Design Principles Applied:}
\begin{itemize}[leftmargin=1.5em]
    \item \textit{Strategy Pattern} — Different LLM models for different agent roles.
    \item \textit{Factory Pattern} — \texttt{get\_llm(purpose)} creates the right model instance.
    \item \textit{Single Responsibility} — Utilities extracted into \texttt{utils.py} (JSON parsing, truncation, memory management, Tavily search).
    \item \textit{Configuration over Convention} — All tunable parameters in \texttt{config.py}.
\end{itemize}
\end{reasonbox}

\begin{impactbox}
This single design session produced 9 concrete improvements that were all implemented in one pass:
\begin{enumerate}
    \item Centralized \texttt{config.py} (no more magic numbers)
    \item Multi-model LLM factory with purpose-based selection
    \item Retry + exponential backoff for rate limits
    \item Automatic model fallback (70B $\to$ 8B on failure)
    \item Robust JSON extraction (handles code fences, partial output)
    \item Sliding window + summary memory management
    \item Optimized Tavily search (top 3, \texttt{include\_answer=True})
    \item LLM-based confidence scoring (replaced naive heuristic)
    \item Dynamic, context-aware clarification questions
\end{enumerate}
\end{impactbox}

%% ═══════════════════════════════════════════════════════════
\section{Prompt 6 — Bug Diagnosis from Production Logs}
%% ═══════════════════════════════════════════════════════════

\begin{promptbox}
\textit{(Pasted the full 413 error traceback from the terminal)}\\[4pt]
\texttt{groq.APIStatusError: Error code: 413 - \{`error': \{`message': `Request too large for model llama-3.1-8b-instant... Limit 6000, Requested 7058'\}\}}
\end{promptbox}

\begin{reasonbox}
When testing revealed a production error, I provided the \textbf{complete error traceback} rather than a vague description. This is a critical engineering practice:

\begin{itemize}[leftmargin=1.5em]
    \item The traceback pinpointed the exact function (\texttt{synthesis\_node}), the exact API call (\texttt{llm.invoke}), and the exact constraint violated (6,000 TPM limit vs. 7,058 requested).
    \item Rather than asking ``why is it broken,'' the prompt structure (raw logs) enabled immediate root-cause analysis.
\end{itemize}

\textbf{Engineering Practice Applied:} \textit{Evidence-Based Debugging} — Always provide full error context; never paraphrase stack traces.
\end{reasonbox}

\begin{impactbox}
This led directly to the token management system: the \texttt{truncate\_text()} utility, the sliding window memory, and the decision to limit conversation history to the last 4 messages in the Synthesis Agent. These changes eliminated all token-overflow errors permanently.
\end{impactbox}

%% ═══════════════════════════════════════════════════════════
\section{Prompt 7 — Enforcing Design Patterns \& Iterative Refinement}
%% ═══════════════════════════════════════════════════════════

\begin{promptbox}
\textit{``The app is functioning now — update the README.md file accordingly and tell me in chat about what is left and then we will work on testing and debugging then documentation part and making the project directory structure well organized and low coupling with high cohesion and following the software design patterns and proper industry grade standards.''}
\end{promptbox}

\begin{reasonbox}
This prompt is significant because it was given \textbf{after the app was already working}. Most developers stop here — ``it works, ship it.'' Instead, I explicitly demanded a structured, multi-phase refinement cycle:

\begin{enumerate}[leftmargin=1.5em]
    \item \textbf{First: Update documentation} — Ensure the README reflects reality before moving on. This prevents documentation drift, a common problem in industry projects.
    \item \textbf{Second: Test and debug} — The app working on one query does not mean it's robust. I wanted systematic testing before any refactoring.
    \item \textbf{Third: Refactor for design patterns} — Only after tests pass do we restructure for maintainability. This follows the \textbf{Red-Green-Refactor} cycle from TDD methodology.
    \item \textbf{Fourth: Documentation polish} — Final documentation reflects the final, clean codebase — not an intermediate state.
\end{enumerate}

The explicit mention of \textbf{``low coupling with high cohesion''} and \textbf{``software design patterns''} shows a deliberate engineering vocabulary. These aren't buzzwords — they directly influenced the refactoring:

\begin{itemize}[leftmargin=1.5em]
    \item \textbf{Low Coupling:} Agents in \texttt{nodes.py} have no knowledge of the HTTP layer in \texttt{main.py}. The graph in \texttt{graph.py} doesn't know about specific prompts. The frontend doesn't know about LangGraph internals.
    \item \textbf{High Cohesion:} All LLM-related utilities (retry, fallback, JSON parsing) live together in \texttt{utils.py}. All configuration lives in \texttt{config.py}. Each agent function handles exactly one responsibility.
    \item \textbf{Factory Pattern:} \texttt{get\_llm(purpose)} encapsulates model selection behind a clean interface.
    \item \textbf{Strategy Pattern:} Different models are assigned to different agents based on task complexity.
\end{itemize}

\textbf{Engineering Principle Applied:} \textit{``Working'' is not ``done.''} — Production-grade code requires deliberate refactoring after functionality is proven.
\end{reasonbox}

\begin{impactbox}
This prompt initiated the largest single improvement pass in the project. It produced:
\begin{itemize}[leftmargin=1.5em]
    \item The extraction of \texttt{src/config.py} — eliminating all magic numbers and hardcoded model names.
    \item The creation of \texttt{src/agent/utils.py} — consolidating 6 shared utilities (LLM factory, retry logic, JSON parsing, truncation, memory management, Tavily search) that were previously scattered across \texttt{nodes.py}.
    \item The README increment tracking table — a professional changelog that demonstrates disciplined, incremental delivery.
\end{itemize}

The codebase went from ``it works'' to ``it's maintainable, extensible, and production-ready'' — which is exactly the quality bar the assignment specified.
\end{impactbox}

%% ═══════════════════════════════════════════════════════════
\section{Prompt 8 — Stress Testing \& Iterative Prompt Refinement}
%% ═══════════════════════════════════════════════════════════

\begin{promptbox}
\textit{``Give me some test cases on which I test it — don't do it by yourself, I'll test it myself.''} \\[4pt]
\textit{(After first round):} ``This is our chat so far — provide critique for the testing.'' \\[4pt]
\textit{(After fixes):} ``We have successfully stress tested our application on multiple prompts — review testing.txt and apply any final touches to further improve our system prompts.''
\end{promptbox}

\begin{reasonbox}
This was not a single prompt but a \textbf{multi-round testing loop} that I drove across the entire development process. The methodology mirrors how QA works in production engineering:

\begin{enumerate}[leftmargin=1.5em]
    \item \textbf{Round 1 — Structured Test Plan:} I requested organized test cases grouped by feature: Clarity Agent, Research Pipeline, Multi-Turn Memory, Edge Cases, and Session Reset. I ran these \textit{manually} rather than having the AI auto-test, ensuring I understood the system's behavior firsthand.
    \item \textbf{Round 1 Results — Bugs Found:}
    \begin{itemize}[leftmargin=1em]
        \item \textbf{Stale State Bug:} When switching topics (Tesla $\to$ Apple), old research data from Tesla polluted the Apple search, causing the Synthesis Agent to report ``no Apple data found.'' Root cause: LangGraph's \texttt{MemorySaver} persisted \texttt{research\_data} across turns.
        \item \textbf{Token Overflow:} Large Tavily results exceeded Groq's 6,000 TPM free-tier limit, crashing the Synthesis Agent with a 413 error.
    \end{itemize}
    \item \textbf{Round 2 — Regression Testing:} After fixes, I re-ran all 10 test cases to verify the fixes worked and no regressions were introduced. Key improvements:
    \begin{itemize}[leftmargin=1em]
        \item Apple topic switch: \textbf{Fixed} — went from broken clarification to clean 8/10 confidence.
        \item Revenue comparison: Confidence jumped from \textbf{5/10 to 9/10} after calibrating the scoring rubric.
        \item Context bleed eliminated: Apple response no longer mentioned Microsoft.
        \item Zero crashes on massive queries (WW2 overview).
    \end{itemize}
    \item \textbf{Prompt Engineering Insight:} During testing, I noticed the Clarity Agent prompts were becoming overly specific — essentially ``hardcoding'' edge cases like \texttt{``if the name is X, AI, Go...''} I identified this as \textbf{overfitting to test data} and refactored all prompts to be \textit{principle-based} rather than example-based. Instead of listing specific entity names, the prompt now teaches the LLM a generalizable decision framework.
\end{enumerate}

\textbf{Engineering Principle Applied:} \textit{Test $\to$ Fix $\to$ Retest $\to$ Generalize.} The testing loop didn't just find bugs — it drove architectural improvements and revealed a meta-lesson about prompt engineering (teach principles, not examples).
\end{reasonbox}

\begin{impactbox}
The two-round stress testing process produced the following concrete improvements:
\begin{itemize}[leftmargin=1.5em]
    \item \textbf{State reset in Clarity Agent:} Added \texttt{research\_data: ""} and \texttt{attempts: 0} to the Clarity node's return value, ensuring every new query starts with a clean research slate.
    \item \textbf{Calibrated confidence rubric:} Replaced the vague ``consider relevance'' prompt with a 5-tier scoring rubric (0-2, 3-4, 5-6, 7-8, 9-10), reducing the 5/10 default clustering.
    \item \textbf{Topic-aware query builder:} Added the rule ``if the latest query names a NEW entity, search for THAT entity only'' to the Research Agent, preventing cross-topic pollution.
    \item \textbf{Principle-based Clarity prompt:} Replaced 4 hardcoded rules with a generalizable decision framework: ``A query is clear if a reasonable person could identify the research target.''
    \item \textbf{Focus directive in Synthesis:} Added ``Focus ONLY on the topic the user is currently asking about'' to eliminate context bleed from prior conversation turns.
\end{itemize}

The full test transcript (\texttt{testing.txt}) is included in the repository as evidence of this systematic quality assurance process.
\end{impactbox}

%% ═══════════════════════════════════════════════════════════
\section{Prompt 9 — Enterprise Readiness \& Real-Time UX}
%% ═══════════════════════════════════════════════════════════

\begin{promptbox}
\textit{``I was thinking to switch the frontend to a proper frontend should I do it or not like React also tell me the backend is good and any suggestions... I have a lot of time left like I just made all this in nearly 2 hours... yes add all''}
\end{promptbox}

\begin{reasonbox}
With the core system stable and tested, I had a choice: rewrite the frontend in React, or optimize the existing stack to industry standards. I chose the latter, deciding that a lightweight vanilla JS frontend paired with a highly optimized backend is a stronger demonstration of core engineering principles than adding framework overhead.

To elevate the project to true \textbf{Enterprise Readiness}, I directed the implementation of four critical additions:
\begin{enumerate}[leftmargin=1.5em]
    \item \textbf{Server-Sent Events (SSE) Streaming:} The standard request/response cycle left the UI unresponsive during the 10-15 seconds of multi-agent processing. By implementing SSE (`FastAPI StreamingResponse`), the backend now pushes real-time agent activation events to the frontend, transforming a "black box" wait into an engaging pipeline animation.
    \item \textbf{Unit Testing (Pytest):} I requested a comprehensive test suite to prove the codebase's modularity. 23 unit tests were written to validate configuration, LLM instantiation, JSON extraction, and API schemas. 
    \item \textbf{Structured Logging:} Replaced all raw `print()` statements with a robust `logger.info/warning/error` module, which is mandatory for production observability.
    \item \textbf{Dockerization:} Added a `Dockerfile` to ensure the application is easily containerized and deployed.
\end{enumerate}

\textbf{Engineering Principle Applied:} \textit{UX is not just React.} Real-time feedback mechanisms (streaming) and backend reliability (tests, logs) provide more value in an AI application than the choice of frontend framework.
\end{reasonbox}

\begin{impactbox}
These final additions bridge the gap between a "student project" and a "production prototype." The SSE implementation, in particular, is a hallmark of modern AI products (like ChatGPT or Perplexity), demonstrating a deep understanding of asynchronous communication and user experience in generative AI systems.
\end{impactbox}

%% ═══════════════════════════════════════════════════════════
\section*{Summary of Problem-Solving Methodology}
%% ═══════════════════════════════════════════════════════════

\begin{center}
\begin{tabular}{|l|l|}
\hline
\textbf{Phase} & \textbf{Approach} \\
\hline
Architecture & Plan before code; define state schema and routing first \\
Technology & Evaluate tools against project constraints, not defaults \\
Structure & Layered architecture, modular directories, living docs \\
Quality & Manual test plan covering happy path, edge cases, memory \\
Optimization & Multi-model strategy, memory management, search tuning \\
Debugging & Full traceback analysis, evidence-based root-cause fixes \\
Refinement & Enforce design patterns after functionality is proven \\
Stress Testing & Multi-round test-fix-retest loop with principle-based prompt refinement \\
\hline
\end{tabular}
\end{center}

\vspace{10pt}

\section*{Key Takeaway}

My prompting strategy mirrors how I approach real-world software engineering: \textbf{think, then build, then test, then improve}. Each prompt was deliberately crafted to produce architectural decisions, not just code output. The AI was used as a \textbf{pair programmer} — I drove the design direction, proposed the improvements, and validated the results. The code is the artifact; the \textit{thinking} is what makes it production-grade.

\end{document}
